\clearpage
\fancyhead[L]{\small COMET.}
\hypertarget{native-464-comet}{}
\begin{center}{\large COMET.}\end{center}

\begin{center}
\rule{0.55\linewidth}{0.4pt}\\[2pt]
Astronomische Nachrichten Band VI. Nr. 123. Seite 43\,..\,44. 1827 September.\\[2pt]
\rule{0.55\linewidth}{0.4pt}
\end{center}

\begin{center}Beobachtungen des von \textsc{Pons} im Luchs entdeckten Cometen.\end{center}

Herr Hofr. \textsc{Gauss} hat mir folgende von Ihm gemachte Meridianbeobachtungen gesandt.

\begin{center}
\begin{tabular}{l|l|l|l}
 & & AR. Com. & $\delta$ Com. \\
\hline
1827 Aug. 20 & $10^{\rm u}28^{\rm m}11^{\rm s}1$ & $125^\circ29'41''9$ & $+56^\circ21'56''1$ \\
\hfill 21 & $37\;53.0$ & $128\;54\;42.5$ & $54\;58\;38.2$ \\
\hfill 22 & $47\;27.0$ & $132\;17\;54.6$ & $53\;23\;1.1$ \\
\end{tabular}
\end{center}

Aus diesen Beobachtungen hat Herr \textsc{Peters} folgende Bahn berechnet:

\Needspace{9\baselineskip}
\SourceLeaf{465}
Zeit des Prihels 1827 September $12.2632$ Altona.

\begin{center}
\begin{tabular}{l l l}
$P$ & \mbox{.\,.\,.\,.\,.\,.\,.\,.\,.\,.\,.\,.\,.\,.\,.\,.\,.\,.\,.} & $255^\circ25'50''$ \\
$\AscendingNode$ & \mbox{.\,.\,.\,.\,.\,.\,.\,.\,.\,.\,.\,.\,.\,.\,.\,.\,.\,.\,.} & $150\;11\;40$ \\
$i$ & \mbox{.\,.\,.\,.\,.\,.\,.\,.\,.\,.\,.\,.\,.\,.\,.\,.\,.\,.\,.} & $54\;27\;5\mathrm{c}$ \\
$\log q$ & \mbox{.\,.\,.\,.\,.\,.\,.\,.\,.\,.\,.\,.\,.\,.\,.\,.\,.\,.} & $9.21395$ \\
\end{tabular}
\end{center}

Rückläufig.

Diese stellen die mittlere Beobachtung auf $+5''$ in Länge und $-4''$ in Breite dar.
\begin{flushright}\textsc{Schumacher.}\end{flushright}

\clearpage
\fancyhead[L]{\small PALLAS.}
\hypertarget{native-465-pallas}{}
\begin{center}{\large PALLAS.}\end{center}

\begin{center}
\rule{0.55\linewidth}{0.4pt}\\[2pt]
Astronomische Nachrichten Band VI. Nr. 124 Seite 67\,..\,68. 1827 October.\\[2pt]
\rule{0.55\linewidth}{0.4pt}
\end{center}

\begin{center}Pallasbeobachtungen.\end{center}

Herr Hofrath \textsc{Gauss} hat mir folgende von Ihm gemachte Pallasbeobachtungen mitgetheilt.

\begin{center}
\begin{tabular}{l|l|l|l}
 & & AR. & $\delta$ \\
\hline
1827 Aug. 31 & $11^{\rm u}16^{\rm m}2^{\rm s}0$ & $328^\circ19'54''6$ & $+5^\circ52'27''4$ \\
\hfill Sept. 2 & $11\;6\;42.2$ & $327\;57\;51.4$ & $5\;28\;38.2$ \\
\hfill 3 & $11\;2\;2.8$ & $327\;46\;57.4$ & $5\;16\;31.8$ \\
\hfill 4 & $10\;57\;24.2$ & $327\;36\;14.9$ & $5\;4\;21.3$ \\
\end{tabular}
\end{center}

Die erste Beobachtung ist weniger zuverlässig, da wegen der grossen Ungleichheit der Temperatur aussen und innen der Planet äusserst schwer zu sehen war; bei den 3 andern waren die Fenster lange vorher geöffnet, und Herr Hofrath \textsc{Gauss} hält diese Beobachtungen für sehr gut.
\begin{flushright}\textsc{Schumacher.}\end{flushright}

\clearpage
\fancyhead[L]{\small GRÖSSTE SONNENHÖHE.}
\hypertarget{native-465-sunheight}{}
\begin{center}{\large GRÖSSTE SONNENHÖHE.}\end{center}

\begin{center}
\rule{0.55\linewidth}{0.4pt}\\[2pt]
Astronomische Hülfstafeln herausg. von \textsc{Schumacher} für 1827 Seite 92.\\[2pt]
\rule{0.55\linewidth}{0.4pt}
\end{center}

Bekanntlich erreicht die Sonne wegen ihrer Veränderung in Declination nicht im Meridian ihre grösste Höhe, sondern in den aufsteigenden Zeichen \emph{nach}, in den niedersteigenden Zeichen \emph{vor} ihrem Durchgange durch den Meridian. Der Augenblick dieser grössten Höhe wird dadurch merkwürdig, dass er nach Herrn Hofraths \textsc{Gauss} sinnreicher Bemerkung, bei Beobachtungen von Circummeridianhöhen
\SourceLeaf{466}der Sonne dem Berechner die Mühe erspart, die Veränderung der Declination während der Beobachtungszeit in Rechnung zu ziehen. Er braucht dazu nur die Stundenwinkel, die als Argumente der Correctionen dienen, durch die man ausser dem Mittage gemessene Höhen auf die Mittagshöhe reducirt, nicht vom wahren Mittage wie gewöhnlich, sondern von dem Augenblicke zu rechnen, in dem die Sonne ihre grösste Höhe erreicht, und hat dann, durch diese einfache Veränderung des Nullpunktes der Stundenwinkel, den Einfluss der Veränderung der Declination für die ganze Ausdehnung seiner Beobachtungsreihe in Rechnung gezogen. --- Die folgende Tafel enthält für jeden zehnten Tag, und für alle Polhöhen von Gibraltar bis Petersburg (von $36^\circ$ bis $60^\circ$) das, was man zu dem Augenblicke des wahren Mittags hinzuzufügen, oder davon abzuziehen hat, um den Augenblick zu erhalten, in dem die Sonne die grösste Höhe erreicht. \textsc{Schumacher.}
